\appendix
\section{Appendix}
\subsubsection{Reinforcement Learning Overview}
Reinforcement learning (RL) models sequential decision making as a Markov decision process (MDP) $(\mathcal{S}, \mathcal{A}, p, r, \gamma)$. At each time step $t$, the agent observes a state $s_t \in \mathcal{S}$, selects an action $a_t \in \mathcal{A}$ according to a policy $\pi_\theta(a_t \mid s_t)$, receives a scalar reward $r_t = r(s_t, a_t)$, and the environment transitions to $s_{t+1} \sim p(\cdot \mid s_t, a_t)$. The objective is to find policy parameters $\theta$ that maximize the expected discounted return
\[
J(\theta) = \mathbb{E}_{\pi_\theta} \Bigg[ \sum_{t=0}^{\infty} \gamma^t r_t \Bigg],
\]
where $\gamma \in (0,1]$ is the discount factor. For a fixed policy $\pi$, the state-value function and action-value function are
\[
V^\pi(s) = \mathbb{E}_\pi \big[ \sum_{t=0}^{\infty} \gamma^t r_t \mid s_0 = s \big], 
\qquad
Q^\pi(s,a) = \mathbb{E}_\pi \big[ \sum_{t=0}^{\infty} \gamma^t r_t \mid s_0 = s, a_0 = a \big],
\]
and they satisfy the Bellman expectation equation
\[
V^\pi(s) = \sum_{a} \pi(a \mid s) \sum_{s'} p(s' \mid s, a) \big[ r(s,a) + \gamma V^\pi(s') \big].
\]
The advantage function is defined as
\[
A^\pi(s,a) = Q^\pi(s,a) - V^\pi(s),
\]
and measures how much better action $a$ is compared to the average action at state $s$.

\paragraph{Policy gradient principle.}
For differentiable policies $\pi_\theta(a \mid s)$, the policy gradient theorem gives
\[
\nabla_\theta J(\theta) = \mathbb{E}_{s \sim d^\pi, a \sim \pi_\theta} \big[ \nabla_\theta \log \pi_\theta(a \mid s) \, Q^\pi(s,a) \big],
\]
and in practice we replace $Q^\pi(s,a)$ by an advantage estimate $\hat A_t$, leading to the classic REINFORCE-style update
\[
\theta \leftarrow \theta + \alpha \, \hat A_t \, \nabla_\theta \log \pi_\theta(a_t \mid s_t).
\]
Most modern RL methods (PPO, GRPO, and their variants) can be viewed as safer or more sample-efficient ways of performing this update.

\paragraph{Proximal Policy Optimization (PPO).}
PPO \cite{schulman2017proximalpolicyoptimizationalgorithms} stabilizes policy gradients by preventing excessively large policy updates. Let
\[
r_t(\theta) = \frac{\pi_\theta(a_t \mid s_t)}{\pi_{\theta_{\text{old}}}(a_t \mid s_t)}
\]
be the probability ratio between the new and the behavior policy. Let $\hat A_t$ be an advantage estimate (e.g., GAE). PPO maximizes the clipped surrogate objective
\[
L^{\text{PPO}}(\theta) = \mathbb{E}_t \bigg[ \min\big( r_t(\theta) \hat A_t,\; \operatorname{clip}(r_t(\theta), 1 - \epsilon, 1 + \epsilon) \hat A_t \big) \bigg],
\]
where $\epsilon$ is a small hyperparameter (e.g., $0.1$ or $0.2$). The $\min(\cdot)$ and the clipping together act as a trust region in probability space, so that even if $\hat A_t$ is large, the policy will not move too far away from $\pi_{\theta_{\text{old}}}$ in a single update.

\paragraph{Group Relative Policy Optimization (GRPO).}
Recent LLM-oriented RL work observes that when many responses can be sampled for the same prompt, we can estimate advantages without a separate value network. GRPO forms a group $\mathcal{G} = \{r_1, \dots, r_N\}$ of rollouts for the same input, scores them with a reward model $R_\phi(\cdot)$, and then normalizes within the group:
\[
A_i = \frac{R_\phi(r_i) - \operatorname{mean}(\mathcal{G})}{\operatorname{std}(\mathcal{G})}.
\]
This gives a critic-free advantage estimate. GRPO then uses essentially the same clipped PPO-style loss
\[
L^{\text{GRPO}}(\theta) = \frac{1}{N} \sum_{i=1}^{N}
\min \bigg(
\frac{\pi_\theta(r_i \mid x)}{\pi_{\theta_{\text{old}}}(r_i \mid x)} A_i,
\operatorname{clip}\Big(\frac{\pi_\theta(r_i \mid x)}{\pi_{\theta_{\text{old}}}(r_i \mid x)}, 1 - \epsilon, 1 + \epsilon \Big) A_i
\bigg)
 - \lambda \, D_{\mathrm{KL}}\big(\pi_\theta \,\|\, \pi_{\text{ref}}\big),
\]
where the KL term keeps the updated policy close to a reference model (often the SFT model). Compared to PPO, the key differences are: (i) advantages come from groupwise normalization, not a learned value function; (ii) the method is well suited to large-batch LLM training where many samples per prompt are available.

\paragraph{Group-in-Group Policy Optimization (GiGPO).}
Very recent work on training LLM agents in multi-step environments argues that a single group-level relative advantage (as in GRPO) is not enough, because the reward can be sparse and credit assignment needs to happen both at the trajectory level and at individual steps. GiGPO \cite{feng2025groupingrouppolicyoptimizationllm} introduces a two-level grouping:
\begin{itemize}
  \item at the episode level, sample multiple full trajectories for the same task and compute a macro relative advantage (similar to GRPO);
  \item at the step level, retroactively group actions that occurred in the same environment state (called anchor states) across different trajectories, and compute a micro relative advantage for those actions.
\end{itemize}
This yields fine-grained credit signals without introducing a value network. Conceptually, GiGPO keeps the good parts of GRPO (critic-free, low memory, PPO-like clipping) but adds an extra axis of grouping so that actions that happened at the same state can be compared directly. This is particularly attractive for long-horizon urban/traffic tasks or LLM agents, where the effect of a single bad action might only become visible many steps later.

\subsubsection{Traffic Signal Control}

\paragraph{General overview.}
Traffic signal control (TSC) aims to coordinate phase sequences and durations at signalized intersections to maximize throughput and reduce congestion. Classical fixed-time methods use the Webster formula\cite{webster} to preset signal timings based on historical traffic data collected from the traffic network at different
times. In contrast, the max-pressure control framework\cite{MERCADER2020275} computes a pressure value for each signal phase from the difference between queue lengths on opposing approaches and selects the phase with the highest pressure; this strategy can be shown to maximize network throughput. Reinforcement-learning (RL) methods model TSC as a Markov decision process in which an agent observes traffic states (for example, lane queue lengths and delays), chooses an action (phase switch) and receives a reward linked to throughput or delay reduction. CoLight\cite{10.1145/3357384.3357902} learns network-level cooperation by letting intersections communicate through graph attention, so signals can coordinate beyond immediate neighbors. CoSLight\cite{ruan2024coslightcooptimizingcollaboratorselection} extends this idea by learning not only the signal policy but also which intersections should collaborate at each moment, improving multi-intersection coordination under changing congestion patterns. MPLight\cite{Chen_Wei_Xu_Zheng_Yang_Xiong_Xu_Li_2020} is a decentralized, pressure-aware RL variant built to run efficiently in large CityFlow\cite{tang2019cityflowcityscalebenchmarkmultitarget}-style networks, using parameter sharing to keep city-scale control tractable. Despite achieving competitive performance in simulation, many RL-based controllers suffer from poor cross-city generalization and lack interpretability, motivating recent research into interpretable and generalizable TSC agents.

\paragraph{CityFlow environment.}
A key enabler for modern RL research on TSC is the availability of scalable simulators. SUMO\cite{SUMO2018}, the traditional open-source traffic simulator, is often too slow for city-scale RL experiments. CityFlow\cite{tang2019cityflowcityscalebenchmarkmultitarget} is a multi-agent microscopic traffic simulator designed for large-scale RL studies. It provides flexible definitions of road networks and traffic flows from both synthetic and real-world data and exposes a user-friendly API for interaction with RL agents. CityFlow's multithreaded design makes it more than twenty times faster than SUMO, enabling efficient simulation of networks with hundreds or thousands of intersections. Because of its speed and extensibility, CityFlow has become a standard environment for evaluating RL-based TSC algorithms.

\paragraph{\textit{LLMLight}: Large language models as TSC agents.}
The \textit{LLMLight}\cite{lai2024llmlightlargelanguagemodels} framework treats traffic signal control as a partially observable Markov game and leverages large language models (LLMs) to serve as decision-making agents. At each signal switch, an agent collects observation features such as the number of queuing vehicles and approaching vehicles on each incoming lane, augments this information with a traffic scenario description, a task description and a commonsense knowledge prompt that encourages the agent to prioritize long queues, and feeds the combined text to an LLM. The LLM performs chain-of-thought reasoning to decide which phase to activate and provides a human-readable rationale. \textit{LLMLight} constructs a specialized backbone model called LightGPT via imitation fine-tuning on GPT-4 trajectories and critic-guided policy refinement. During imitation, the model minimizes the negative log-likelihood
\begin{equation*}
\mathcal{L}_{\mathrm{NLL}} = - \sum_{t=1}^{T} \log p_\theta \bigl( y_t \mid y_{<t}, \pi, \mathcal{P} \bigr),
\end{equation*}
where $\pi$ is the prompt consisting of observation and commonsense instructions and $\mathcal{P}$ denotes the reasoning trajectory. A critic model derived from an action-value network evaluates sampled trajectories, and a ranking-based loss steers the model toward trajectories with higher long-term rewards. Experimental results show that \textit{LLMLight} with LightGPT achieves competitive or superior performance and provides better interpretability compared with traditional RL baselines across diverse traffic scenarios. It provides a reproducible workflow that integrates CityFlow with RL-based fine-tuning of LLMs (though implemented in TensorFlow) and datasets ready for use.

\paragraph{\textit{Traffic-R1}: Reinforced LLMs for human-like TSC.}
\textit{Traffic-R1}\cite{zou2025trafficr1reinforcedllmsbring} is a 3B foundation model trained to perform traffic signal control with human-like reasoning. It adopts a two-stage reinforcement-learning fine-tuning paradigm: an offline stage aligns the model with question-answer (QA) samples (generated by DeepSeek-R1 and verified by human experts), and an online stage allows self-exploration in simulation to refine the policy via reward feedback. The reward is calculated from the cumulative queuing time after multiple steps, and each step only half of the agents are activated, whose decisions are propagated to neighboring agents as an input. This mimics an asynchronous communication network allows multiple agents to coordinate across intersections. Benchmark results report that \textit{Traffic-R1} offers zero-shot generalization to unseen road networks and out-of-distribution incidents, outperforming both classical heuristics and RL controllers. When deployed in production, \textit{Traffic-R1} manages signals affecting more than 55,000 drivers each day, reducing average queue lengths by over 5\% and significantly lowering operator workload. Though without a codebase, \textit{Traffic-R1} provides insights into using SOTA RL algorithms like GRPO to train TSC agents, as well as the innovative idea of mimicking the real world scenario with a asynchronous mult-agent setup.

\subsubsection{Joint-Scored GRPO Comparison Variant}
\label{appendix:joint_scored_grpo}
This appendix summarizes the \textit{Joint-Scored GRPO} pipeline used as a comparison variant to the proposed \textit{QGuided-GRPO} method. The purpose of this variant is to test whether indirect short-horizon joint return estimation can provide a better training signal than direct critic guidance. Unlike the main method, it does not rely on a pretrained action-value critic to score each sampled response directly.

\paragraph{Snapshot-forked joint candidate evaluation.}
At each decision step, the policy samples $k$ candidate responses for every intersection and parses them into actions. Instead of scoring each sampled action independently, the method aligns candidates across intersections into $k$ joint action vectors and evaluates each vector by forking the current simulator snapshot. Each fork starts from the same environment state but executes a different sampled joint action, allowing the method to compare candidate action sets under a common traffic condition.

\paragraph{Short-horizon discounted return.}
For each forked candidate $c$, the simulator is rolled out for a short horizon $H$. After the first jointly sampled action is executed, subsequent actions are resampled from the current policy at later steps. The candidate score is the discounted return
\[
G_t^{(c)} = \sum_{\tau=0}^{H-1} \gamma^\tau r_{t+\tau}^{(c)},
\]
where $\gamma$ is the discount factor and $r_{t+\tau}^{(c)}$ is the congestion-based network reward observed at rollout step $t+\tau$ in candidate fork $c$.

\paragraph{Local, neighborhood, and global reward mixing.}
The per-step reward is computed from a mixed congestion cost that combines local intersection congestion, neighboring congestion, and global network congestion. Let $m_i$ denote the local congestion metric at intersection $i$ (for example, incoming queue length or pressure), let $\bar m_{\mathcal{N}(i)}$ denote the mean metric over the neighbors of $i$, and let $\bar m_{\mathcal{V}}$ denote the mean metric over the whole network. The reward is defined from the negative mixed cost
\[
r_t = - \frac{1}{|\mathcal{V}|} \sum_{i \in \mathcal{V}}
\Bigl(
\alpha m_i
  + \beta \bar m_{\mathcal{N}(i)}
  + (1-\alpha-\beta)\bar m_{\mathcal{V}}
\Bigr),
\]
where $\alpha$ and $\beta$ control the relative contribution of local and neighborhood congestion. This construction encourages decisions that reduce both immediate local congestion and broader network imbalance.

\paragraph{Projection back to per-intersection action scores.}
Because GRPO training still expects prompt-level rewards for individual sampled responses, the joint candidate returns must be projected back to action-level scores. For each intersection $i$ and each action $a \in A$, let $\mathcal{C}_{i,a}$ be the set of joint candidates whose parsed action at intersection $i$ equals $a$. The induced action score is defined as the mean return of those candidates:
\[
\tilde Q_{i,t}(a) =
\frac{1}{|\mathcal{C}_{i,a}|}
\sum_{c \in \mathcal{C}_{i,a}} G_t^{(c)},
\]
with a fallback reward used when no sampled candidate proposes action $a$. The resulting vector $\tilde{\mathbf{q}}_{i,t}$ plays the same role in GRPO training as the critic score vector $\mathbf{q}_{i,t}$ in the main method, but it is obtained indirectly through joint rollout evaluation.

\paragraph{Collect-then-train split.}
In implementation, the variant is organized as a collect-then-train pipeline. A rollout stage first saves prompt records, selected actions, and projected action-score vectors obtained from snapshot-forked candidate evaluation. The GRPO stage then reloads these records and fine-tunes the LLM policy using the same prompt-conditioned reward lookup mechanism as in the main method. This design isolates the additional cost of joint candidate evaluation from the policy optimization stage and makes it suitable for controlled comparison against \textit{QGuided-GRPO}.


\subsubsection{Sample Prompt}
\label{appendix:sample_prompt}
You are an expert in traffic management.

A traffic light regulates a four-way intersection with northern, southern, eastern, and western approaches, each containing two lanes: one for through traffic and one for left-turns. Each lane is further divided into three segments. Segment 1 is the closest to the intersection. Segment 2 is in the middle. Segment 3 is the farthest. In a lane, there may be early queued vehicles and approaching vehicles traveling in different segments. Early queued vehicles have already arrived at the intersection and await passage permission. Approaching vehicles will arrive at the intersection in the future.

The traffic light has 4 signal phases. Each signal relieves vehicles' flow in a group of two specific lanes.

Available signal phases:
- ETWT: Eastern and western through lanes.
- NTST: Northern and southern through lanes.
- ELWL: Eastern and western left-turn lanes.
- NLSL: Northern and southern left-turn lanes.

Current intersection state:
Signal: ETWT
Relieves: Eastern and western through lanes.
- Early queued: 0 (East), 1 (West), 1 (Total)
- Segment 1: 0 (East), 0 (West), 0 (Total)
- Segment 2: 1 (East), 0 (West), 1 (Total)
- Segment 3: 0 (East), 2 (West), 2 (Total)
- Neighbor incoming totals: 2 (East), 1 (West), 3 (Known total), 2/2 available

Signal: NTST
Relieves: Northern and southern through lanes.
- Early queued: 0 (North), 0 (South), 0 (Total)
- Segment 1: 0 (North), 0 (South), 0 (Total)
- Segment 2: 0 (North), 0 (South), 0 (Total)
- Segment 3: 1 (North), 0 (South), 1 (Total)
- Neighbor incoming totals: NA (North), 3 (South), 3 (Known total), 1/2 available

Signal: ELWL
Relieves: Eastern and western left-turn lanes.
- Early queued: 0 (East), 0 (West), 0 (Total)
- Segment 1: 0 (East), 0 (West), 0 (Total)
- Segment 2: 0 (East), 0 (West), 0 (Total)
- Segment 3: 0 (East), 1 (West), 1 (Total)
- Neighbor incoming totals: 2 (East), 1 (West), 3 (Known total), 2/2 available

Signal: NLSL
Relieves: Northern and southern left-turn lanes.
- Early queued: 0 (North), 0 (South), 0 (Total)
- Segment 1: 0 (North), 0 (South), 0 (Total)
- Segment 2: 0 (North), 0 (South), 0 (Total)
- Segment 3: 0 (North), 0 (South), 0 (Total)
- Neighbor incoming totals: NA (North), 3 (South), 3 (Known total), 1/2 available

The state description above lists:
- The group of lanes relieved under each traffic light phase.
- The number of early queued vehicles in the allowed lanes of each signal.
- The number of approaching vehicles in different segments of the allowed lanes of each signal.
- Neighbor incoming totals from adjacent intersections for each phase.
- `NA` means that adjacent side has a virtual/missing neighbor and is excluded from `Known total`.

Question:
Which is the most effective traffic signal that will most significantly improve the traffic condition during the next phase?

Note:
- Traffic congestion is primarily dictated by early queued vehicles, with the most significant impact.
- You must pay the most attention to lanes with long queue lengths.
- It is not urgent to consider vehicles in distant segments, since they are unlikely to reach the intersection soon.

Requirements:
- Think step by step.
- You can only choose one of the signals listed above.
- Step 1: Provide a brief analysis identifying the optimal traffic signal.
- Step 2: After finishing the analysis, answer with your chosen signal.
- Include exactly one final decision tag in this format: <signal>PHASE</signal>, where PHASE is one of: ETWT, NTST, ELWL, NLSL.